\documentclass[10pt,a4paper]{amsart}
\usepackage[margin=19mm]{geometry}
\usepackage{amsmath,amssymb,mathtools}
\usepackage[scr=rsfs,cal=euler]{mathalfa}
\usepackage{xfrac}
\usepackage[unicode,hidelinks,bookmarks=false]{hyperref}
\newcommand{\ie}{i.e.\ }
\newcommand{\cf}{cf.\ }
\newcommand{\resp}{respectively\ }
\newcommand{\Subsection}{\subsection}
\providecommand{\nobreakdash}{\nobreak\hbox{-}}
\setlength{\emergencystretch}{2em}
\multlinegap0pt
\usepackage{amssymb}
\usepackage{mathtools}
\usepackage[shortlabels]{enumitem}
\usepackage{bm}
\usepackage{tensor}
\usepackage{float}
\newtheorem{proposition}{Proposition}
\newcommand{\pol}{\mathrm{P}}
\title{Spectral Properties of\\ Nekrasov--Okounkov Polynomials}
\author{Bernhard Heim, Markus Neuhauser}
\date{Source: arXiv:2606.15394v2 (CC BY 4.0)}
\begin{document}
\maketitle

\noindent\emph{Source and licence.} Spectral Properties of\\ Nekrasov--Okounkov Polynomials. Version arXiv:2606.15394v2 (CC BY 4.0), distributed by arXiv under the Creative Commons Attribution 4.0 International licence, http://creativecommons.org/licenses/by/4.0/, as recorded in the arXiv licence metadata for this exact version and shown on the abstract page https://arxiv.org/abs/2606.15394v2. Full licence notice: \texttt{licenses/arxiv-2606.15394-CC-BY-4.0.txt}.\par\smallskip

\noindent\textit{Editorial scope.} Counted unit: the article's proposition Hessenberg (source0.tex lines 713--717). The article's complete original proof of it is delivered with it (source0.tex lines 719--815, one \texttt{proof} environment of 97 lines). No mathematics is written, repaired or re-derived: the statement and the proof are the article's own lines, in the article's own notation, and no retained line is substituted or re-flowed.  No further source-local context is needed: every cross-reference of every retained line resolves inside the two delivered spans.  Nothing was omitted from any retained span, so the package carries no omission record.  No retained line contains a citation, so the wrapper carries no bibliography and no bibliography entry.  No retained line contains a figure environment, an image-inclusion command or drawing code, so no image is delivered and the package declares no asset dependency.  Editorial formatting only, in addition: an AMS \texttt{amsart} preamble that restates the article's own declarations and macros used by the retained lines, namely the environments \texttt{proposition} on separate counters, together with the standard packages those lines need.  The retained environments print in this document as Proposition 1 in order of appearance, whereas the article prints the same items with the numbering of its own full document.  The complete native source of the article as distributed in the arXiv e-print for this exact version is bundled unchanged as \texttt{sources/202854/source0.tex}.
\begin{proposition} We have
\begin{equation*} \label{Hessenberg}
\pol _n^{h}(z) =\det \left( z H_n^{-1} L_n - U_n \right).
\end{equation*}
\end{proposition}

\begin{proof}
The proof is by mathematical induction on $n$.
The case $n=1$ is obvious as $\pol _{1}^{h}\left( z\right) =z$.
Suppose now $n\geq 2$.
We use
the induction hypothesis that
$\pol _{j}^{h}\left( z\right) =\det \left( zH_{j}^{-1}L_{j}-U_{j}\right) $
for all $j<n$.

The matrix $zH_{n}^{-1}L_{n}-U_{n}$ has the following form as determinant
\[
\left|
\begin{array}{ccccc}
zh\left( 1\right) ^{-1}\sigma \left( 1\right) &-1
&0
&\cdots &0\\
zh\left( 2\right) ^{-1}\sigma \left( 2\right)
&zh\left( 2\right) ^{-1}\sigma \left( 1\right) &-1
&\ddots &0\\
\vdots &\vdots
&\ddots &\ddots &\vdots \\
zh\left( n-1\right) ^{-1}\sigma \left( n-1\right)
&zh\left( n-1\right) ^{-1}\sigma \left( n-2\right) &\cdots
&zh\left( n-1\right) ^{-1}\sigma \left( 1\right) &-1\\
zh\left( n\right) ^{-1}\sigma \left( n\right)
&zh\left( n\right) ^{-1}\sigma \left( n-1\right) &\cdots
&zh\left( n\right) ^{-1}\sigma \left( 2\right)
&zh\left( n\right) ^{-1}\sigma \left( 1\right)
\end{array}
\right| .
\]
We expand this with respect to the last row. So we have to compute the
determinant of the matrices obtained by
deleting
the last row and
the $k$th
column. This results in a matrix determinant
\[
\left|
\begin{array}{ccccc}
zh\left( 1\right) ^{-1}\sigma \left( 1\right) &-1
&0&\cdots
&0\\
zh\left( 2\right) ^{-1}\sigma \left( 2\right)
&zh\left( 2\right) ^{-1}\sigma \left( 1\right)
&-1&\ddots
&0\\
\vdots &\vdots
&\ddots
&\ddots &\vdots \\
zh\left( n-2\right) ^{-1}\sigma \left( n-2\right)
&zh\left( n-2\right) ^{-1}\sigma \left( n-3\right) &\cdots
&-1&0\\
zh\left( n-1\right) ^{-1}\sigma \left( n-1\right)
&zh\left( n-1\right) ^{-1}\sigma \left( n-2\right) &\cdots
&zh\left( n-1\right) ^{-1}\sigma \left( 1\right) &-1
\end{array}
\right| .
\]
Since we deleted column $k$
the
$\left( k-1\right) \times \left( n-k\right) $ block in the upper right corner
contains only zeros, i.~e.\ in rows $1$ to $k-1$ and (original) columns $k+1$ to
$n$. Therefore, we can compute the determinant as the product of the upper
left block of size $\left( k-1\right) \times \left( k-1\right) $ and the
lower right block of size $\left( n-k\right) \times \left( n-k\right) $.
Note that this lower right
block is a lower triangular matrix with $-1$ on the
diagonal so the determinant is $\left( -1\right) ^{n-k}$. The upper left
block is exactly the matrix $zH_{k-1}^{-1}L_{k-1}-U_{k-1}$. Therefore, by
expansion of the determinant
\begin{eqnarray*}
&&\det \left( zH_{n}^{-1}L_{n}-U_{n}\right) \\
&=&%\left( -1\right) ^{n+1}
\frac{z}{h\left( n\right)
}\sigma \left( n\right)
%\left( -1\right) ^{n-1}%\\
%%&&{}
+\sum _{k=2}^{n}
%%\left( -1\right) ^{n+k}
\frac{z}{h\left( n\right)
}\sigma \left(
n+1-k\right) \det \left( zH_{k-1}^{-1}L_{k-1}-U_{k-1}\right) 
%%\left( -1\right)^{n-k}
\\
&=&\frac{z}{h\left( n\right)
}\sigma \left( n\right) +\sum _{k=2}^{n}\frac{z}{h\left(
n\right)
}\sigma \left( n+1-k\right) \pol _{k-1}^{h}\left( z\right) \\
&=&\frac{z}{h\left( n\right) }\left( \sigma \left( n\right) +\sum _{j
=1}^{n-1}\sigma \left( j
\right) \pol _{n-j
}^{h}\left( z\right) \right)
=
\pol _{n}^{h}\left( z\right) .
\end{eqnarray*}
\end{proof}

\end{document}
